\section{A Lagrangian comparison along a closed geodesic}
\label{sec:lagrangian}

We work in the intermediate case
\[
 D=\frac{\pi}{2},\qquad T=2D=\pi,\qquad
 d=n-1,\qquad 0<k<d,\qquad m=d-k.
\]
Fix a unit-speed geodesic $\gamma:\R\to M$.  The structural conclusions
in Proposition~\ref{geom:structure} give the common period $T$, the absence of
conjugate points at separations in $(0,D)$, and conjugate multiplicity
$k$ at separation $D$.  This section uses these geometric conclusions
and the Jacobi equation.  Its comparison requires no assumption about
the volume of $M$.

The purpose of the construction is to separate a Jacobi determinant
into two factors with different scalar comparison profiles.  Unlike
the curvature-invariant splitting supplied by the totally geodesic
curvature-four sphere families of
\citet[Theorem~26]{RovenskiiToponogov1996}, this separation takes place
in a fixed coordinate chart on the space of Jacobi fields.  No
splitting of the curvature operator is assumed.

\subsection{The symplectic curve of vanishing Jacobi fields}

We use the classical symplectic description of Jacobi fields;
see, for example, \citet[Sections~2.1.2 and~10.5.2]{Zelditch2017Notes}
and the variational-index framework of \citet{Duistermaat1976}.
We derive the required identities to fix the signs and the
normalization of the endpoint pairing.

Let $\mathcal J$ be the real vector space of normal Jacobi fields along
$\gamma$.  A prime on a field denotes covariant differentiation along
$\gamma$.  Prescribing the normal pair $(J(t),J'(t))$ at one time
identifies $\mathcal J$ with a vector space of dimension $2d$.  The form
\begin{equation}
 \omega(I,J)=\langle I,J'\rangle-\langle I',J\rangle
 \label{eq:lag-wronskian}
\end{equation}
is independent of the evaluation time: its derivative vanishes by the
Jacobi equation and the self-adjointness of
$X\mapsto R(X,\dot\gamma)\dot\gamma$.  Evaluation of the initial data
identifies it with the standard nondegenerate symplectic form.

For each $t\in\R$, set
\[
 L_t=\{J\in\mathcal J:J(t)=0\}.
\]
The Jacobi initial-value problem shows that $L_t$ depends smoothly on
$t$, has dimension $d$, and is isotropic for $\omega$.  It is therefore
Lagrangian.  Its motion carries a natural positive inner product.
For $I,J\in L_t$, choose a smooth local section $J_s\in L_s$ with
$J_t=J$, and define
\[
 G_t(I,J)=\omega(I,\dot J_t),
 \qquad
 \dot J_t=\left.\frac{d}{ds}\right|_{s=t}J_s\in\mathcal J.
\]
The definition is independent of the extension: a section of $L_s$
vanishing as an element of $\mathcal J$ at $s=t$ has derivative in
$L_t$, which pairs to zero with $I$.  Differentiating $J_s(s)=0$ gives
$\dot J_t(t)=-J'(t)$.  Since $I(t)=0$, our sign convention yields
\begin{equation}
 G_t(I,J)=\langle I'(t),J'(t)\rangle.
 \label{eq:lag-velocity}
\end{equation}
Thus $G_t$ is symmetric and positive definite, and
$J\mapsto J'(t)$ is an isometry from $(L_t,G_t)$ to the normal space
$\dot\gamma(t)^\perp$.

We will need periodicity of the fields themselves.  For any normal
initial data $(a,b)$ at time zero, the connection splitting of
$T(SM)$ realizes $(a,b)$ as the derivative of a smooth curve of unit
initial velocities.  The associated variation through unit-speed
geodesics has Jacobi field with initial data $(a,b)$; it is normal
because its tangential component and its derivative vanish initially.
All geodesics in this variation satisfy the same identity
$\phi_T=\Id_{SM}$.  Differentiating that identity, or differentiating
the equality of the varied positions and velocities after time $T$,
gives
\begin{equation}
 J(t+T)=J(t),\qquad J'(t+T)=J'(t)
 \quad(J\in\mathcal J).
 \label{eq:lag-periodicity}
\end{equation}
These are equalities in the same tangent space, since
$\gamma(t+T)=\gamma(t)$.  They do not assert that parallel transport
around $\gamma$ is the identity.  In particular,
\[
 (L_{t+T},G_{t+T})=(L_t,G_t).
\]

The intersection pattern is consequently
\begin{equation}
 \dim(L_x\cap L_y)=
 \begin{cases}
  0,&0<y-x<T,\quad y-x\ne D,\\
  k,&y-x=D.
 \end{cases}
 \label{eq:lag-intersections}
\end{equation}
For $0<y-x<D$, this is exponential nonsingularity.  At separation
$D$, evaluation at $y=x+D$ of the fields vanishing at $x$ identifies the
kernel with the kernel of the angular differential of the cut-locus
fibration, whose dimension is $k$.  If $D<y-x<T$, a field in the
intersection also vanishes at $x+T$ by
\eqref{eq:lag-periodicity}; the separation from $y$ to $x+T$ is less
than $D$, so the field is zero.

\subsection{The determinant and a fixed graph chart}

\begin{proposition}[Jacobi determinant normalization]
\label{lag:determinant}
Choose $G_x$- and $G_y$-orthonormal frames
$Z_x=(Z_{x,1},\ldots,Z_{x,d})$ and
$Z_y=(Z_{y,1},\ldots,Z_{y,d})$ of $L_x$ and $L_y$.  Then
\begin{equation}
 f(x,y)=\left|\det\bigl(\omega(Z_{x,a},Z_{y,b})\bigr)_{a,b=1}^d\right|
 \label{eq:lag-f}
\end{equation}
is independent of these choices.  It is continuous, symmetric in its
arguments, and separately $T$-periodic.  It is the absolute normal
Jacobi determinant at $y$ of fields vanishing at $x$ with orthonormal
initial derivatives there.  In particular, for $0<y-x<D$ it is the
polar volume density with respect to $dt$ and ordinary sphere area at
$\gamma(x)$.

Use the same orthonormal basis of $\dot\gamma(0)^\perp$ for the
coordinates
\[
 (q,p)=(J(0),J'(0))\in\R^d\oplus\R^d,
\]
choosing its first $k$ vectors to be the initial derivatives of a basis
of $L_0\cap L_D$.  On
\[
 \mathcal I=(0,D)\cup(D,T)
\]
there is a smooth symmetric matrix $A(t)$ such that
$L_t=\{(q,A(t)q):q\in\R^d\}$, and $A'(t)>0$.  For $x,y\in\mathcal I$,
\begin{equation}
 f(x,y)=
 \frac{|\det(A(y)-A(x))|}
 {\sqrt{\det A'(x)\det A'(y)}}.
 \label{eq:lag-graph-determinant}
\end{equation}
\end{proposition}

\begin{proof}
Changing either orthonormal frame multiplies the pairing matrix by an
orthogonal matrix.  This leaves its absolute determinant unchanged.
Local smooth orthonormal frames exist, so $f$ is continuous even when
the pairing is singular.  Skew-symmetry of $\omega$ replaces the matrix
by its negative transpose when the arguments are exchanged, proving
symmetry.  Periodicity follows from
\eqref{eq:lag-periodicity} and \eqref{eq:lag-velocity}.

Evaluation of the pairing at $y$ gives
\[
 \omega(Z_{x,a},Z_{y,b})
 =\langle Z_{x,a}(y),Z'_{y,b}(y)\rangle.
\]
The vectors $Z'_{y,b}(y)$ form an orthonormal normal basis, and the
fields $Z_{x,a}$ have zero values and orthonormal derivatives at $x$.
This proves the claimed determinant interpretation.  More explicitly,
an angular variation of a unit direction with derivative $e$ gives the
Jacobi field $J$ with $J(x)=0$ and $J'(x)=e$.  Its value at $x+r$ is the
angular derivative of $\exp_{\gamma(x)}(r\,\cdot)$.  Together with the
unit radial derivative and Gauss's lemma, this gives the polar density
$f(x,x+r)\,dr\,d\sigma$.  There is no additional power of $r$ in this
polar formula; the Cartesian normal-coordinate density would be
$f(x,x+r)/r^d$.

In the stated coordinates,
\[
 \omega((q,p),(\widetilde q,\widetilde p))
 =q^{\mathsf T}\widetilde p-p^{\mathsf T}\widetilde q,
 \qquad L_0=\{q=0\}.
\]
Equation~\eqref{eq:lag-intersections} makes $L_t$ transverse to $L_0$
for $t\in\mathcal I$, proving the graph representation.  Isotropy is
equivalent to $A=A^{\mathsf T}$.  For the graph frame
$E_tq=(q,A(t)q)$,
\[
 G_t(E_tq,E_tr)=q^{\mathsf T}A'(t)r,
 \qquad
 \omega(E_xq,E_yr)=q^{\mathsf T}(A(y)-A(x))r.
\]
Positivity of $G_t$ gives $A'(t)>0$.  The frames
$E_tA'(t)^{-1/2}$ are orthonormal for $G_t$, so taking determinants in
the last identity proves \eqref{eq:lag-graph-determinant}.
\end{proof}

The coordinate basis is now fixed.  Split both $q$ and $p$ into their
first $k$ and last $m$ coordinates and write
\begin{equation}
 A(t)=
 \begin{pmatrix}
  V_b(t)&C(t)\\ C(t)^{\mathsf T}&H(t)
 \end{pmatrix},
 \qquad
 S(t)=V_b'(t)-C'(t)H'(t)^{-1}C'(t)^{\mathsf T}.
 \label{eq:lag-blocks}
\end{equation}
The subscript in $V_b$ distinguishes this matrix from the volume of the
manifold.  Here $H$ is an $m\times m$ matrix and $S$ is a $k\times k$
matrix.  Positive definiteness of $A'$ and block elimination give
\begin{equation}
 H'(t)>0,\qquad S(t)>0,\qquad
 \det A'(t)=\det H'(t)\det S(t)
 \quad(t\in\mathcal I).
 \label{eq:lag-velocity-factorization}
\end{equation}

For comparison, on the standard model the adapted chart has
\[
 A_{\mathrm{mod}}(t)
 =\diag\bigl(-2\cot(2t)\Id_k,-\cot(t)\Id_m\bigr).
\]
Thus $H'_{\mathrm{mod}}(t)=\csc^2(t)\Id_m$ and
$S_{\mathrm{mod}}(t)=4\csc^2(2t)\Id_k$.
These formulas motivate the scalar reference profiles below.

\subsection{The chart poles and determinant-root profiles}

The curve $L_t$ remains smooth at $0,D,T$; it is the projection from
$L_t$ to the fixed $q$-coordinate space that ceases to be invertible.
Its kernel has dimension $d$ at $0,T$ and dimension $k$ at $D$.
The following expansion shows how this loss of transversality enters
$A$: at $D$ the pole lies entirely in the first $k$ coordinates,
while the complementary block $H$ remains regular.  Smooth remainders
are needed because the comparison uses derivatives of the graph matrices.

\begin{lemma}[Pole expansions]
\label{lag:poles}
There is a positive definite $k\times k$ matrix $B$ such that the graph
matrix has the expansions
\begin{align}
 A(t)&=-\frac1{t-D}
 \begin{pmatrix}B&0\\0&0\end{pmatrix}+R_D(t),
 &&\text{near }D,\label{eq:lag-middle-pole}\\
 A(t)&=-\frac1{t-t_0}\Id_d+R_{t_0}(t),
 &&\text{near }t_0=0,T,\label{eq:lag-end-poles}
\end{align}
where each remainder extends smoothly across its indicated point.
In particular, $C$ and $H$ extend smoothly across $D$, and
$H'(D)>0$.
\end{lemma}

\begin{proof}
Put $\tau=t-D$.  Choose a local smooth frame of $L_t$, written as the
columns of
\[
 F(t)=\begin{pmatrix}X(t)\\P(t)\end{pmatrix},
\]
whose first $k$ columns at $D$ are the vertical vectors $(0,e_a)$ for
$1\le a\le k$.  These are exactly the chosen coordinate basis of
$L_0\cap L_D$.  Isotropy with these vectors forces every vector in the
image of $X(D)$ to have its first $k$ coordinates zero.  The kernel of
the value projection on $L_D$ is $L_0\cap L_D$, so that image has
dimension $m$ and is the entire last coordinate space.  Choose the
remaining frame columns so that their value projections give a basis
of this image.  In blocks, their last $m$ value coordinates form an
invertible matrix $M$.

The first $k$ columns of $X(D)$ vanish.  Dividing these columns by
$\tau$ gives a smooth matrix $\widetilde X(t)$ satisfying
\[
 X(t)=\widetilde X(t)
 \begin{pmatrix}\tau\Id_k&0\\0&\Id_m\end{pmatrix}.
\]
Let $G$ be the matrix of $G_D$ on the first $k$ frame vectors.  For
$1\le a,b\le k$, the definition of the velocity form in the fixed
coordinates gives
\[
 G_{ab}=\omega((0,e_a),\dot F_b(D))
 =-e_a^{\mathsf T}X_b'(D).
\]
Consequently
\begin{equation}
 \widetilde X(D)=
 \begin{pmatrix}-G&0\\N&M\end{pmatrix},
 \qquad G>0,\quad M\text{ invertible},
 \label{eq:lag-divided-frame}
\end{equation}
for some matrix $N$.  In particular $\widetilde X(D)$ is invertible.
For $\tau\ne0$ sufficiently small,
\[
 A(t)=P(t)
 \begin{pmatrix}\tau^{-1}\Id_k&0\\0&\Id_m\end{pmatrix}
 \widetilde X(t)^{-1}.
\]
To identify its residue, let $P_1(t)$ denote the first $k$ columns of
$P(t)$ and let $U_1(t)$ denote the first $k$ rows of
$\widetilde X(t)^{-1}$.  Then the only possibly singular summand is
$\tau^{-1}P_1(t)U_1(t)$.  At $D$,
\[
 P_1(D)=\begin{pmatrix}\Id_k\\0\end{pmatrix},
 \qquad U_1(D)=\begin{pmatrix}-G^{-1}&0\end{pmatrix}.
\]
Their product is $-\diag(G^{-1},0)$.  The difference
$P_1(t)U_1(t)-P_1(D)U_1(D)$ divided by $\tau$ is smooth, as is the
summand from the other frame columns.  This proves
\eqref{eq:lag-middle-pole} with $B=G^{-1}>0$, including the absence of
any off-diagonal residue.

For fixed $z\in\R^m$, the graph vectors
\[
 E_z(t)=\left(\begin{pmatrix}0\\z\end{pmatrix},
                 \begin{pmatrix}C(t)z\\H(t)z\end{pmatrix}\right)
\]
extend smoothly through $D$.  Their limits lie in $L_D$ by continuity
of $L_t$, and are nonzero if $z\ne0$.  The velocity identity in the
fixed symplectic coordinates, now also at $D$, gives
\[
 z^{\mathsf T}H'(D)z
 =\omega(E_z(D),\dot E_z(D))
 =G_D(E_z(D),E_z(D))>0.
\]
This proves strict positivity of the extended derivative.

At $t_0=0$, take a frame whose value at $0$ is the whole vertical
coordinate basis.  All $X$-columns then vanish at $0$.  The velocity
metric on this basis is $\Id_d$ by \eqref{eq:lag-velocity}; hence
$X'(0)=-\Id_d$ and $P(0)=\Id_d$.  Dividing all value columns by $t$
and repeating the preceding calculation gives residue $-\Id_d$.
At $t_0=T$, periodicity identifies both $L_T$ and its velocity metric
with their values at zero in these same fixed coordinates.  The same
argument gives residue $-\Id_d$ there, proving
\eqref{eq:lag-end-poles}.
\end{proof}

\begin{proposition}[Scalar profiles]
\label{lag:profiles}
Define
\begin{equation}
 \begin{aligned}
 h(t)&=\det(H'(t))^{1/m} &&(0<t<T),\\
 s(t)&=\det(S(t))^{1/k} &&(t\in\mathcal I),\\
 s_i(t)&=s(t+iD) &&(0<t<D,\ i=0,1).
 \end{aligned}
 \label{eq:lag-profiles}
\end{equation}
The function $h$ is smooth and positive across $D$.  There are positive
constants $c_H,C_H,c_i,C_i$, depending on the fixed geodesic and its
chart, such that
\begin{align}
 c_H&\le h(t)\sin^2t\le C_H &&(0<t<T),
 \label{eq:lag-h-bounds}\\
 c_i&\le \frac{s_i(t)\sin^2(2t)}4\le C_i
 &&(0<t<D,\ i=0,1).
 \label{eq:lag-s-bounds}
\end{align}
Equivalently, the logarithms of the ratios of $h$ to $1/\sin^2t$ and
of $s_i$ to $4/\sin^2(2t)$ are bounded.
\end{proposition}

\begin{proof}
Lemma~\ref{lag:poles} and \eqref{eq:lag-velocity-factorization} show
that $H'$ is positive throughout $(0,T)$, including $D$.  Near
$t_0=0,T$, put $\tau=t-t_0$.  Differentiating the smooth-remainder
expansion gives
\[
 H'(t)=\tau^{-2}\Id_m+O(1),\qquad
 V_b'(t)=\tau^{-2}\Id_k+O(1),\qquad C'(t)=O(1).
\]
Here and below $O(1)$ denotes a matrix bounded in operator norm as the
indicated time is approached.  Since $H'(t)^{-1}=O(\tau^2)$, the Schur
complement satisfies
\begin{equation}
 S(t)=\tau^{-2}\Id_k+O(1)\qquad(t\to0,T).
 \label{eq:lag-S-end}
\end{equation}
At $D$, both $C'$ and $H'$ are smooth and $H'(D)$ is invertible.  The
middle pole therefore yields
\begin{equation}
 S(t)=(t-D)^{-2}B+O(1)\qquad(t\to D).
 \label{eq:lag-S-middle}
\end{equation}
Taking determinant roots gives the positive limits
\[
 \begin{array}{c|cc}
 &\text{left endpoint}&\text{right endpoint}\\ \hline
 h(t)\sin^2t\quad(0<t<T)&1&1\\
 s_0(t)\sin^2(2t)/4\quad(0<t<D)&1&\det(B)^{1/k}\\
 s_1(t)\sin^2(2t)/4\quad(0<t<D)&\det(B)^{1/k}&1.
 \end{array}
\]
Each ratio thus extends continuously and positively to its compact
closed interval.  Its minimum and maximum supply the stated bounds.
\end{proof}

\subsection{Schur complements for arbitrary endpoint order}

The Schur complement and determinant-root concavity are standard
matrix tools \citep[Appendix~A.5.5 and Exercises~3.18(b), 3.26(b)]
{BoydVandenberghe2004}. Here the essential point is to apply the
factorization with either actual endpoint order. We retain the
matrix calculation so that its signs and integration paths are explicit.

For distinct $x,y\in\mathcal I$, use $\Delta$ to denote the difference
of the values at $y$ and $x$.  Since $H$ is smooth and $H'>0$ throughout
$(0,T)$,
\begin{equation}
 \Delta H=H(y)-H(x)=\int_x^y H'(t)\,dt
 \label{eq:lag-H-difference}
\end{equation}
is positive definite for $y>x$ and negative definite for $y<x$.
In particular it is invertible in either order.  Define
\begin{equation}
 Q(x,y)=\Delta V_b-(\Delta C)(\Delta H)^{-1}(\Delta C)^{\mathsf T}.
 \label{eq:lag-Q}
\end{equation}
Thus $Q$ is the Schur complement of the regular block in the endpoint
difference $\Delta A$, whereas $S$ in \eqref{eq:lag-blocks} is the
Schur complement in the velocity $A'$.  The next proposition relates
them by differentiating in either endpoint; the diagonal and cut-pair
zeros of $Q$ then anchor the two interval bounds.

\begin{proposition}[Two Schur-complement bounds]
\label{lag:schur}
The matrix $Q(x,y)$ is symmetric and satisfies
\begin{equation}
 \det(\Delta A)=\det(\Delta H)\det Q(x,y),\qquad
 Q(y,x)=-Q(x,y).
 \label{eq:lag-Schur-factorization}
\end{equation}
It extends continuously by zero to $y=x\in\mathcal I$, and
\begin{equation}
 Q(u,u+D)=Q(u+D,u)=0\qquad(0<u<D).
 \label{eq:lag-Q-zeros}
\end{equation}
Whenever its distinct endpoints remain in $\mathcal I$, it obeys the
positive-semidefinite matrix inequalities
\begin{equation}
 \partial_yQ(x,y)\ge S(y),\qquad
 -\partial_xQ(x,y)\ge S(x).
 \label{eq:lag-Q-derivatives}
\end{equation}
Consequently, for every $0<u<v<D$ and $i,j\in\{0,1\}$,
\begin{align}
 Q(u+iD,v+jD)&\ge\int_u^v S(t+jD)\,dt,
 \label{eq:lag-Q-bound-j}\\
 Q(u+iD,v+jD)&\ge\int_u^v S(t+iD)\,dt.
 \label{eq:lag-Q-bound-i}
\end{align}
In particular, this $Q$ is positive definite for all four choices of
$i,j$, including $i=1,j=0$.
\end{proposition}

\begin{proof}
With $R=(\Delta C)(\Delta H)^{-1}$, block congruence gives
\begin{equation}
 \begin{pmatrix}\Id_k&-R\\0&\Id_m\end{pmatrix}
 (\Delta A)
 \begin{pmatrix}\Id_k&0\\-R^{\mathsf T}&\Id_m\end{pmatrix}
 =\begin{pmatrix}Q&0\\0&\Delta H\end{pmatrix}.
 \label{eq:lag-block-congruence}
\end{equation}
This proves the determinant identity and identifies the nullity of
$\Delta A$ with that of $Q$.  Exchanging the endpoints changes the
sign of each difference, and hence changes $Q$ to $-Q$.

For a matched pair $(u,u+D)$, the kernel of $\Delta A$ parametrizes
$L_u\cap L_{u+D}$ by the common $q$ coordinate.  It has dimension $k$
by \eqref{eq:lag-intersections}.  The matrix $Q$ has size $k\times k$
and therefore has full nullity: it is the zero matrix.  This proves
\eqref{eq:lag-Q-zeros}, also in the reversed order.  Near a diagonal
point within the chart, Taylor expansion gives
\[
 \Delta H=(y-x)H'(x)+O((y-x)^2),\quad
 \Delta C=(y-x)C'(x)+O((y-x)^2),
\]
and the corresponding expansion for $\Delta V_b$.  Invertibility of
$H'(x)$ then gives
\begin{equation}
 Q(x,y)=(y-x)S(x)+O((y-x)^2).
 \label{eq:lag-Q-diagonal}
\end{equation}
This establishes the zero extension needed at the endpoints of the
integrations below.

Direct differentiation of \eqref{eq:lag-Q}, using
$\partial_y(\Delta H)^{-1}=-(\Delta H)^{-1}H'(y)(\Delta H)^{-1}$,
gives
\begin{align*}
 \partial_yQ
 &=V_b'(y)-C'(y)R^{\mathsf T}-R C'(y)^{\mathsf T}
                  +R H'(y)R^{\mathsf T}\\
 &=\begin{pmatrix}\Id_k&-R\end{pmatrix}A'(y)
                  \begin{pmatrix}\Id_k\\-R^{\mathsf T}\end{pmatrix}.
\end{align*}
The identical calculation with $-\partial_x$ gives the same expression
with $A'(x)$.  For either endpoint $a=x$ or $a=y$, completing the
square gives the exact identity
\begin{align}
 &\begin{pmatrix}\Id_k&-R\end{pmatrix}A'(a)
                  \begin{pmatrix}\Id_k\\-R^{\mathsf T}\end{pmatrix}
 \notag\\
 &\quad=S(a)+
 \bigl(R-C'(a)H'(a)^{-1}\bigr)H'(a)
 \bigl(R-C'(a)H'(a)^{-1}\bigr)^{\mathsf T}.
 \label{eq:lag-square-completion}
\end{align}
The second summand is positive semidefinite.  This proves
\eqref{eq:lag-Q-derivatives}.  The calculation used invertibility of
$\Delta H$, but did not require its positive definiteness.

To obtain \eqref{eq:lag-Q-bound-j}, fix the first endpoint $u+iD$ and
move the second from $u+jD$ to $v+jD$.  The starting value of $Q$ is
zero by the diagonal extension if $i=j$, and by
\eqref{eq:lag-Q-zeros} otherwise.  Integrating $\partial_yQ\ge S(y)$
on this interval gives the result.  For
\eqref{eq:lag-Q-bound-i}, fix the second endpoint $v+jD$ and move the
first from $u+iD$ to $v+iD$.  The terminal value is zero, so integration
of $-\partial_xQ\ge S(x)$ gives the second result.  When $i=j$, one
can first integrate on a shortened interval and pass to the diagonal
endpoint using \eqref{eq:lag-Q-diagonal}.

Figure~\ref{fig:schur-endpoint-paths} displays these integrations
with the half-period indices held fixed.
\begin{figure}[tbp]
\centering
\begin{tikzpicture}[x=1.05cm,y=1.05cm,>=Stealth,font=\small]
  \draw[->] (0,0)--(4.7,0) node[right] {$a$};
  \draw[->] (0,0)--(0,4.6) node[above] {$b$};
  \draw[gray,densely dashed] (0,0)--(4,4);
  \draw[gray,dotted] (4,0)--(4,4)--(0,4);
  \node[below left] at (0,0) {$0$};
  \node[below] at (4,0) {$D$};
  \node[left] at (0,4) {$D$};
  \foreach \a/\b in {1/u,3/v} {
    \draw (\a,-.06)--(\a,.06) node[below=3pt] {$\b$};
    \draw (-.06,\a)--(.06,\a) node[left=3pt] {$\b$};
  }
  \draw[gray,dotted] (1,0)--(1,1);
  \draw[gray,dotted] (0,3)--(1,3);
  \draw[very thick,->] (1,1)--(1,3);
  \draw[very thick,->] (1,3)--(3,3);
  \fill (1,1) circle (1.7pt);
  \fill (1,3) circle (1.7pt);
  \fill (3,3) circle (1.7pt);
  \node[below right] at (1,1) {$Q(u+iD,u+jD)=0$};
  \node[anchor=east] (qtarget) at (-.3,3.55)
    {$Q(u+iD,v+jD)$};
  \draw[gray,thin] (qtarget.east)--(.9,3.08);
  \node[right,align=left] at (3,3) {$Q(v+iD,v+jD)=0$};
  \node[rotate=90,above] at (.89,2) {$\partial_yQ\ge S(y)$};
  \node[above] at (2,3.02) {$-\partial_xQ\ge S(x)$};
  \node[align=center] at (2,5.05)
    {$x=a+iD,\quad y=b+jD,\quad i,j\in\{0,1\}$};
  \node[align=center] at (2,-.9)
    {$0<u<v<D$\quad (folded endpoint coordinates)};
\end{tikzpicture}

\caption{The two paths in folded endpoint coordinates
$x=a+iD$, $y=b+jD$, for $0<u<v<D$.
The first starts at $Q=0$ and increases the second endpoint;
the second increases the first endpoint toward another zero.
The zeros come from the diagonal when $i=j$ and from the cut-pair
identity otherwise. Each moving endpoint stays in one open
half-period. When $i=1,j=0$, the actual order is $u+D>v$,
and both positive-semidefinite lower bounds still apply.}
\label{fig:schur-endpoint-paths}
\end{figure}

Both moving endpoints stay in one open half-period, so neither path
passes through a chart pole.  In particular, when $i=1,j=0$ the actual
endpoint order is $u+D>v$.  The first path starts at $(u+D,u)$ and the
second ends at $(v+D,v)$; the same zero values and the same derivative
inequalities apply.  Thus neither matrix inequality reverses in this
case.  Finally, the integrals on their right sides are positive
definite, so $Q(u+iD,v+jD)>0$.
\end{proof}

\subsection{Scalar lower bounds for the two factors}

For distinct $x,y\in\mathcal I$, define
\begin{equation}
 f_H(x,y)=
 \frac{|\det(H(y)-H(x))|}
 {\sqrt{\det H'(x)\det H'(y)}},\qquad
 f_K(x,y)=
 \frac{|\det Q(x,y)|}
 {\sqrt{\det S(x)\det S(y)}}.
 \label{eq:lag-two-factors}
\end{equation}
These auxiliary factors are symmetric under exchanging the actual
endpoints.  Combining \eqref{eq:lag-graph-determinant},
\eqref{eq:lag-velocity-factorization}, and
\eqref{eq:lag-Schur-factorization} gives
\begin{equation}
 f(x,y)=f_H(x,y)f_K(x,y).
 \label{eq:lag-factorization}
\end{equation}
The factors are positive away from the cut pairs $|y-x|=D$; at those
pairs $f_K=0$.

\begin{proposition}[Scalar block bounds]
\label{lag:block-bounds}
With the profiles of Proposition~\ref{lag:profiles}, every pair of
distinct chart endpoints satisfies
\begin{equation}
 f_H(x,y)^{1/m}\ge
 \frac{\left|\int_x^y h(t)\,dt\right|}{\sqrt{h(x)h(y)}}.
 \label{eq:lag-H-scalar-bound}
\end{equation}
For $0<u<v<D$ and $i,j\in\{0,1\}$ one has
\begin{equation}
 f_K(u+iD,v+jD)^{1/k}\ge
 \frac{
  \left(\int_u^v s_i(t)\,dt\right)^{1/2}
  \left(\int_u^v s_j(t)\,dt\right)^{1/2}}
 {\sqrt{s_i(u)s_j(v)}}.
 \label{eq:lag-K-scalar-bound}
\end{equation}
Together with endpoint exchange, these bounds cover the full square
$(0,T)^2$ except the chart, diagonal, and cut lines, a set of planar
measure zero.
\end{proposition}

\begin{proof}
We first record the determinant inequality used in both blocks.  For
a positive definite $r\times r$ matrix $M$, write
$\delta_r(M)=\det(M)^{1/r}$.  The elementary variational formula
\begin{equation}
 \delta_r(M)=
 \inf_{\substack{P>0\\\det P=1}}\frac1r\tr(PM)
 \label{eq:lag-determinant-variational}
\end{equation}
follows by applying the arithmetic--geometric mean inequality to the
eigenvalues of $M^{1/2}PM^{1/2}$.  Equality is attained at
$P=\delta_r(M)M^{-1}$.  The formula proves monotonicity under the
positive-semidefinite order.  It also proves, for any continuous
positive definite matrix function on a compact interval,
\begin{equation}
 \delta_r\left(\int_a^b M(t)\,dt\right)
 \ge\int_a^b\delta_r(M(t))\,dt.
 \label{eq:lag-Minkowski}
\end{equation}
Indeed, for every admissible $P$ the integral of
$\tr(PM(t))/r$ is at least the right side, and one then takes the
infimum over $P$.  This is the integral form of Minkowski's determinant
inequality; it retains the degree-one homogeneity of the determinant
root.

Apply \eqref{eq:lag-Minkowski} to $H'$ on the interval between $x$ and
$y$.  Equation~\eqref{eq:lag-H-difference} gives
\[
 |\det(\Delta H)|^{1/m}
 \ge\left|\int_x^y\det(H'(t))^{1/m}\,dt\right|.
\]
For $y<x$ this uses the positive matrix $-\Delta H$.  Division by
\[
 \bigl(\det H'(x)\det H'(y)\bigr)^{1/(2m)}
 =\sqrt{h(x)h(y)}
\]
proves \eqref{eq:lag-H-scalar-bound}.

Now put $Q=Q(u+iD,v+jD)>0$.  Each of the two matrix bounds in
Proposition~\ref{lag:schur}, followed by determinant monotonicity and
\eqref{eq:lag-Minkowski}, gives
\[
 \delta_k(Q)\ge\int_u^v s_i(t)\,dt,
 \qquad
 \delta_k(Q)\ge\int_u^v s_j(t)\,dt.
\]
Thus $\delta_k(Q)$ dominates their maximum and therefore their
geometric mean.  The denominator in the $k$th root of $f_K$ is exactly
\[
 \bigl(\det S(u+iD)\det S(v+jD)\bigr)^{1/(2k)}
 =\sqrt{s_i(u)s_j(v)},
\]
which proves \eqref{eq:lag-K-scalar-bound} with the displayed powers.

Finally, every endpoint in $\mathcal I$ has a unique expression
$u+iD$ with $0<u<D$ and $i\in\{0,1\}$.  For a pair whose folded
parameters differ, order those parameters as $u<v$, exchanging the
actual endpoints if necessary.  The preceding bounds apply for every
choice of the two half-period indices.  Equal folded parameters mean
either a diagonal pair or a cut pair.  The remaining excluded pairs
have an endpoint at $D$.  These are finitely many lines, and symmetry
of both factors justifies the endpoint exchange everywhere else.
\end{proof}

The two determinant roots have now been compared with scalar interval
integrals.  Proposition~\ref{lag:profiles} supplies exactly their
endpoint control.  The next section combines these bounds by a weighted
logarithmic identity.
